\section{Implementation details}
\label{app:implementation}

\subsection{Code structure}

The C++ simulator lives in \texttt{Code/ogcode/cpp\,conversion/} and is structured around four executables, each sharing a common library of routines:
\begin{itemize}[itemsep=2pt, topsep=2pt]
    \item \texttt{ogcode\_cpp.exe} --- baseline learner (no frictions).
    \item \texttt{ogcode\_latency.exe} --- latency-only learner; reads \texttt{A\_Latency.txt}.
    \item \texttt{ogcode\_noise.exe} --- noise-only learner; reads \texttt{A\_NoiseStd.txt}.
    \item \texttt{ogcode\_async.exe} --- asynchronous-actions learner; reads \texttt{A\_LockIn.txt}.
    \item \texttt{ogcode\_combined.exe} --- unified learner accepting all three friction parameters.
\end{itemize}
The single-friction binaries exist primarily for cross-validation against the unified one. Each defaults its non-applicable friction parameters to zero, so at all-zero inputs every binary reproduces the baseline bit-for-bit. This is the principal correctness check.

The unified \texttt{LearningSimulationCombined.cpp} is the implementation actually used for the factorial. The order of operations within an iteration mirrors the description in Section~\ref{sec:methodology-combined}: (1) compute per-agent perceived state with fresh Gaussian noise on the latency-lagged opponent observation, (2) determine the moving agent under $T$-period lock-in, (3) action selection on the perceived state for the moving agent (non-mover repeats their previous price), (4) compute the true joint reward, (5) Q-update each agent on its own perceived state, (6) advance buffers.

\subsection{Random number generation and reproducibility}

The Calvano reference code uses a Numerical-Recipes \texttt{Ran2}-style generator for exploration draws and Q-initialisation tie-breaking. Our translation preserves that generator and its seeding convention so that at zero-friction inputs the unified learner reproduces the baseline exactly. Gaussian noise (Friction~2) uses an independent \texttt{std::mt19937\_64} stream seeded deterministically from the session index. At $\sigma=0$ the noise stream is short-circuited and never advanced, which guarantees that the baseline is reproduced bit-for-bit when noise is off, regardless of how the noise stream is configured.

\section{Design choices not specified in the proposal}
\label{app:design-choices}

This appendix collects every implementation-level decision that the proposal underspecified or where the implementation deliberately departs from the proposal's wording. Each is justified briefly.

\paragraph{Friction~2 noise distribution.}
The proposal specified a discrete-flip noise model: with probability $q$, the rival's observed price index is perturbed to a neighbouring grid point. The implemented version is a Gaussian-on-index model (Equation~\ref{eq:noise}). The Gaussian model is a one-parameter generalisation: at $\sigma=0$ it is the no-noise baseline, at small $\sigma$ it produces small deviations to nearby indices (so it locally agrees with the discrete-flip intuition), and at large $\sigma$ it converges on near-uniform random observation. The discrete-flip variant remains a future-work item.

\paragraph{Friction~3 alternation schedule.}
The proposal specified ``alternating or random selection'' for which agent moves each period; the implemented version uses strict deterministic alternation with a tunable lock-in length $T$ (Section~\ref{sec:methodology-async}). Strict alternation matches the canonical Maskin--Tirole specification and is invariant to RNG state, simplifying cross-validation. Random alternation remains a future-work item.

\paragraph{Off-path deviation tests.}
The proposal specified off-path deviation tests for the asynchronous case as a way to distinguish genuine collusion from timing artefacts. The implemented version does not include them. They are listed as priority extensions in Section~\ref{sec:lim-method}.

\paragraph{Noise on price index, not on price value.}
Equation~\ref{eq:noise} adds noise to the price index $k_j\in\{1,\dots,m\}$ rather than to the price value $p_j$ in money units. The two differ by a multiplicative factor (the spacing of the grid). The index-space convention has the advantage that $\sigma$ has a transparent ``noise in number of grid steps'' interpretation, and the friction parameter is independent of the choice of price grid. In the price-value convention $\sigma$ would have to be re-scaled when the grid changes.

\paragraph{Per-agent noise draws.}
Each agent draws its own independent Gaussian noise each period. Correlated noise across agents (a common shock to monitoring) is a different story (it overlaps with the literature on common observation errors); we do not pursue it here.

\paragraph{Both agents update Q under asynchrony.}
Under strict alternation only the moving agent revises its action. The implementation has both agents Q-update each period regardless of who moved (Section~\ref{sec:methodology-async}). The non-mover's ``action'' is its previous-period price, which is a valid element of its action set, and the non-mover observes the actual joint reward. Discarding the non-mover's update would discard a legitimate learning signal.

\paragraph{Initial-period buffer initialisation.}
The latency buffer of length $L+1$ per agent is filled at session start by repeating the agent's randomly drawn initial price across all slots, rather than by drawing $L+1$ independent initial prices. Repetition is chosen for simplicity; during the first $L$ iterations of every session the exploration probability $\varepsilon\approx 1$, so the buffer contents are essentially overwritten by random exploration before they affect any learnt policy.

\section{Numerical results in full}
\label{app:numerics}

\subsection{Single-friction sweeps}

Tables~\ref{tab:latency-results}, \ref{tab:noise-results} and \ref{tab:async-results} in Section~\ref{sec:results-singles} are the canonical numerical results for the three single-friction sweeps. All values are means of converged session-level $\Delta$, with cross-session standard deviation reported alongside.

\subsection{Grid-robustness check}

Table~\ref{tab:grid-robustness} reports $\Delta$ at four grid sizes $m\in\{15,25,50,100\}$ on the otherwise-identical Calvano baseline (no frictions). The $m=15$ and $m=25$ runs use $1{,}000$ sessions; $m=50$ and $m=100$ use $250$ sessions to keep runtime feasible on a local workstation. The $m=100$ runs additionally use a 5{,}000-episode convergence cap rather than $1{,}000$, since the larger state space takes longer to converge.

\begin{table}[H]
\centering
\caption{Grid-robustness check. $\Delta$ falls modestly with finer price grid but plateaus by $m=50$ and is essentially unchanged from $m=50$ to $m=100$. All runs use Calvano's original fixed-$\beta$ protocol.}
\label{tab:grid-robustness}
\begin{tabular}{rcccc}
\toprule
$m$ & sessions & $\Delta$ & SD & avg time-to-convergence \\
\midrule
15  & 1000 & 0.848 & 0.109 & 261 episodes  \\
25  & 1000 & 0.806 & 0.076 & 364 episodes  \\
50  & 250  & 0.707 & 0.023 & 476 episodes  \\
100 & 250  & 0.706 & 0.016 & 1144 episodes \\
\bottomrule
\end{tabular}
\end{table}

The $\nu$-corrected version of this check, in which $\beta$ is scaled with $m$ to hold expected per-cell visit count constant per Equation~\ref{eq:nu-visit-count}, is left for future work. Without that correction it is conceivable that some of the $\Delta$ decline from $m=15$ to $m=100$ reflects worse exploration per cell rather than a genuine effect of the finer grid; with the correction the comparison would be cleaner.\todo{Run the $\nu$-corrected version on the cluster as part of the same job set as the pairwise heatmaps. Calvano (2020, Section 6.7) discuss this confound.}
